% F17 fig:noise (Chapter 3, margin): the V of the central difference (Central<2>) on exp(x) at x = 1, clean and with
% the values perturbed by a relative noise eta = 1e-10 and 1e-8 (DS19, data/step-error-noise.csv, the perturbation
% of code/noise.hpp). The rounding branch lifts to eta |f| / h and the minimum moves right and up. The default step
% (\valRootEpsThreeDouble, since s = 1) and the step eta^(1/3) for eta = 1e-8 with its error (\valNoiseEightFit...)
% are marked. Steps 1 to 1e-8 (on a log axis the domain restriction is in natural logarithms: ln(1e-8) = -18.42).
\begin{tikzpicture}
\begin{axis}[eqtufte, width=36mm, height=32mm, xmode=log, ymode=log,
    xlabel={step $h$}, ylabel={relative error},
    xmin=1e-8, xmax=1, ymin=1e-12, ymax=1,
    xtickten={-8,-4,0}, ytickten={-12,-8,-4,0},
    restrict x to domain=-18.43:0.01]
  \addplot[muted, mark=*] table[col sep=comma, x=h, y=clean, comment chars={\#}] {data/step-error-noise.csv};
  \addplot[black!75, mark=*] table[col sep=comma, x=h, y=noise_ten, comment chars={\#}] {data/step-error-noise.csv};
  \addplot[accent, mark=*] table[col sep=comma, x=h, y=noise_eight, comment chars={\#}] {data/step-error-noise.csv};
  \draw[muted, thin] ({axis cs:\valRootEpsThreeDouble,1}|-{rel axis cs:0,0})
    -- ({axis cs:\valRootEpsThreeDouble,1}|-{rel axis cs:0,1})
    node[above, font=\sffamily\scriptsize, inner sep=1.5pt] {default};
  \fill[accent] (axis cs:\valNoiseEightFitStep,\valNoiseEightFitRelError) circle (1.2pt);
  \node[accent, font=\sffamily\scriptsize, anchor=south west, inner sep=1pt] at (axis cs:2e-5, 1e-3)
    {$\eta = 10^{-8}$};
  \node[black!75, font=\sffamily\scriptsize, anchor=north west, inner sep=1pt] at (axis cs:1e-7, 1e-6)
    {$10^{-10}$};
  \node[muted, font=\sffamily\scriptsize, anchor=south west, inner sep=1pt] at (axis cs:1e-7, 1e-9)
    {clean};
\end{axis}
\end{tikzpicture}
